% ============================================================
% File:    job_talk.tex
% Purpose: 30-slide beamer job-talk version of main.tex
% Inputs:  ../../output/reg/*_present.tex, ../../output/sum/*_present.tex,
%          ../../output/fig/*.png, img/*, citation.bib
% Outputs: job_talk.pdf
% Author:  EK  Date: 2026-08-31
% ============================================================

\documentclass{beamer}
\usetheme{Boadilla}
\setbeamertemplate{footline}[page number]
\setbeamertemplate{navigation symbols}{}

\usepackage{graphicx}
\usepackage{booktabs}
\usepackage{array}
\usepackage{adjustbox}
\usepackage{varwidth}
\usepackage{amsmath}
\usepackage{amssymb}

\usepackage[
backend=biber,
style=authoryear,
citestyle=bwl-FU
]{biblatex}
\addbibresource{citation.bib}

% ============================================================
% Pipeline output directory (relative to this file)
%   \presreg{name}  -> ../../output/reg/name_present.tex
%   \pressum{name}  -> ../../output/sum/name_present.tex
%   \presfig{f.png} -> ../../output/fig/f.png
% ============================================================
\newcommand{\outputdir}{../../output}

% The generated *_present.tex files are written as article floats:
% \begin{table}[htbp] ... \caption{...} ... \end{table}. Beamer has no float
% mechanism, and wrapping a float in adjustbox raises
% "! LaTeX Error: Not in outer par mode." Neutralise the float and swallow the
% caption (on a slide the frame title carries the caption text), so the table
% body is plain material that can be boxed and scaled.
\makeatletter
\renewenvironment{table}[1][]{\par\centering}{\par}
\renewcommand{\caption}[2][]{}
\makeatother

% Box the table body at its natural width inside a varwidth, so multi-panel
% files still break between panels instead of setting them side by side, then
% shrink only as much as is needed to fit the slide.
% Optional argument overrides the width key, for source files that hardcode a
% narrow adjustbox (e.g. width = 0.45\textwidth) and so need scaling UP: max
% width only ever shrinks.
\newcommand{\slidetableinput}[2][max width=\textwidth]{%
  \adjustbox{#1, max totalheight=0.80\textheight, center}{%
    \begin{varwidth}{17cm}%
      \input{#2}%
    \end{varwidth}}}
\newcommand{\presreg}[2][max width=\textwidth]{\slidetableinput[#1]{\outputdir/reg/#2_present}}
\newcommand{\pressum}[2][max width=\textwidth]{\slidetableinput[#1]{\outputdir/sum/#2_present}}
\newcommand{\presfig}[2][0.75\textheight]{%
  \begin{center}%
  \adjustbox{max width=\textwidth, max height=#1}{\includegraphics{\outputdir/fig/#2}}%
  \end{center}}

% One-line takeaway under a table or figure. The group must close AFTER \par or
% the paragraph inherits the outer font's baselineskip on its first line.
\newcommand{\takeaway}[1]{\par\medskip{\footnotesize #1\par}}

% In-text link to a backup/appendix slide. #1 = \hypertarget name on the
% backup frame, #2 = short description of what the link shows.
\newcommand{\backuplink}[2]{\hyperlink{#1}{\footnotesize(#2 $\rightarrow$)}}

% Underbrace whose label wraps to the width of the braced math, so long labels
% sit under their own term instead of running across the slide.
% #1 = braced math (typeset in display style), #2 = label text.
\newlength{\ubwidth}
\newcommand{\ubtext}[2]{%
  \settowidth{\ubwidth}{$\displaystyle #1$}%
  \underbrace{#1}_{\parbox[t]{\ubwidth}{\centering\scriptsize #2}}}

\AtBeginSection[]{
  \begin{frame}
  \vfill
  \centering
  \begin{beamercolorbox}[sep=8pt,center,shadow=true,rounded=true]{title}
    \usebeamerfont{title}\insertsectionhead\par%
  \end{beamercolorbox}
  \vfill
  \end{frame}
}

\title[Contamination and Self-Reported Compliance]{The Effect of Contamination on Contamination Limit Regulation}
\subtitle{US Coal Mining and Drinking Water Utilities}
\author{Emil Kee-Tui}
\institute{Charles H. Dyson School of Applied Economics and Management \\ Cornell University}
\date{October 2026}

\begin{document}

% ---------------------------------------------------------------- 1
\begin{frame}
\titlepage
\end{frame}

\begin{frame}{Self-reported compliance}
\begin{itemize}
    \item Regulatory models requiring entities to self-report compliance:
    \begin{itemize}
        \item HIPAA violation disclosure; anti-competitiveness/negligence to the Federal Energy Regulatory Commission; SOX reporting.
    \end{itemize}
    \item Incentives to shirk reporting:
    \begin{itemize}
        \item Hide non-compliance
        \item Save self-reporting costs
    \end{itemize}
    \item Environmental regulation on contamination monitored with self-reporting: $\uparrow$ contamination $\rightarrow$ $\uparrow$ incentives 
\end{itemize}

\begin{center}
\begin{beamercolorbox}[sep=10pt,center,rounded=true,shadow=true]{block body}
\large What happens to firm self-reporting when contamination rises?
\end{beamercolorbox}
\end{center}

\begin{center}
\begin{beamercolorbox}[sep=10pt,center,rounded=true,shadow=true]{block body}
\large Are decreases in self-reporting driven by concealment or self-reporting costs?
\end{beamercolorbox}
\end{center}
\end{frame}

\begin{frame}{Drinking water utility self-reporting and coal mining}
\begin{itemize}
    \item I document the effect of contamination on self-reporting at US drinking water utilities exposed to coal mining water contamination between 1985 and 2005. 
    \item Coal mining contaminates ground and surface water through acid mine drainage.
    \item Estimate utility self-reporting, regulator visits, and enforcement on an instrument of upstream coal production. 
    \item $\uparrow$ upstream mining 
    \begin{itemize}
        \item $\uparrow$ arsenic, selenium, and barium concentrations
        \item $\Delta$ in recorded contaminant limit violations = 0
        \item $\downarrow$ self-reporting of water quality falls by 3--5 pp (BASE RATE)
        \item $\uparrow$ \textbf{inspection visits} ($+$2.1 pp) (BASE RATE).
        \item $\downarrow$ Formal enforcement: $-$1.1 pp (base rate of 3\%) (not significant) .
    \end{itemize}
    \item \textbf{Self-reporting costs, not concealment}: 
    \begin{itemize}
        \item When nitrate readings cross 50\% of the MCL $\rightarrow$ monitoring requirements $\uparrow$ (reporting cost $\uparrow$) $\rightarrow$ self-reporting $\downarrow$.
        \item Sanitary inspections at under-reporting utilities do not discover concealed violations.
        \item Contaminant concentrations stay below MCL.
    \end{itemize}
    \end{itemize}
\end{frame}

\begin{frame}{Utility self-reporting nitrate contamination}
    \begin{itemize}
    \item US drinking water utilities:
        \begin{itemize}
            \item Source from ground and/or surface water.
            \item Publicly or privately owned.
            \item Serve twenty-five to millions of customers.
            \item Self-report levels of 83 contaminants by 1992 (cognitive burden).   
        \end{itemize}
    
    \item Self-reporting requirement for nitrates:
    \begin{itemize}
        \item \textbf{Baseline}: carefully collect and promptly send one sample from every intake once a year to a lab
        \item \textbf{Elevated}: N mg/L $\geq$ 5 mg/L, test quarterly
    \end{itemize}

    \item Missing any test is a monitoring and reporting (``MR'') violation.
    \item Violate and regulator conducts a visit
    \item Or enforcement, while remedial rather than punitive, is still costly:
\begin{itemize}
    \item Informal (letters, notices, compliance plans)
    \item Formal (consent decrees, fines, prosecution)
    \end{itemize} 

    \item \cite{oleckno1982national} burden despite low financial cost (National monitoring costs for all IOCs were 2.5 million dollars a year in 1991.) 
  
    \end{itemize}
\end{frame}

\begin{frame}{Data}
\begin{itemize}\small
    \item \textbf{Violations, enforcement, visits} --- utility $\times$ year binary, universe of utilities 1985--2005.
    \item 1479 arsenic, nitrate, and ioc MR violations compared to 49 mcl violations
    \item \textbf{Contaminant concentrations for arsenic, nitrate, barium, selenium} --- 1998--2005.
    \item \textbf{Mine level production and location}
    \item \textbf{Coal sulfur \%} --- of each watershed.
    \item \textbf{Watershed connectivity}.     
    \item Average distance between centroid of adjacent watersheds is 11 km.
    \item \textbf{Sample}: Utilities with a coal mine \textbf{within two watersheds upstream} of their intake and no coal mine colocated with their intake. 10,641 utility-years, 565 utilities.
\end{itemize}
\end{frame}

\begin{frame}{Design: cumulative production, measured concentrations}
\begin{equation*}
\textit{Concen}_{ct}= \beta\cdot\!\!\sum_{\tau=1985}^{t}\!\!\text{UpstreamCoalProd}_{c\tau} + \rho_c + HUC02_c\cdot\tau_t + \epsilon_{ct}
\end{equation*}
\begin{itemize}\small
    \item $\textit{Concen}_{ct}$: mean annual concentration at utility $c$, year $t$ (SYR2, 1998--2005).
    \item $\beta$: effect of 10 million additional cumulative short tons mined upstream since 1985.
    \item $\rho_c$: utility fixed effects. $HUC02_c\cdot\tau_t$: HUC02-specific year effects, absorbing watershed-wide industrial trends.
    \item SEs clustered at the utility level.
\end{itemize}
\medskip
\begin{itemize}
    \item Continuous treatment, \textbf{no control}: identification comes from relative exposure.
    \item Assumption: utilities with large upstream production changes would have trended like those with small ones. \backuplink{balancetest}{balance test on utility characteristics}
\end{itemize}
\end{frame}

% ---------------------------------------------------------------- 18
\begin{frame}{Mining raises contaminant levels}
\presreg{6yr_huc02fe_inorg_ravalli_2005_k2}
\takeaway{Nitrate $+$0.09 mg/L per 10M tons (***): 11.0\% of its mean. Selenium $+$0.0004 mg/L (**): 8.5\% of its mean. \textbf{Every reading in the sample stays under its MCL.} \backuplink{kgridsyr2}{effect across choice of upstream watershed radius}}
\end{frame}

\begin{frame}{The Acid Rain Program (``ARP") moved coal production for reasons unrelated to water quality}
\begin{columns}[T]
\begin{column}{0.44\textwidth}
\begin{itemize}\small

\item Adapt \textcite{stratford2015coalprod} ARP instrument for coal production to
address endogeneity of mining and utility characteristics.
\item 1990 Clean Air Act amendments cap SO$_2$ from power plants; Phase I deadline \textbf{1995}.
\item Power plants were 90\% of coal demand; switching to low sulfur coal was the main compliance response \parencite{lattanzio2022clean}.
\item High sulfur mines have limited ability to wash the sulfur out (bound to chemical structure).
\end{itemize}
\end{column}
\begin{column}{0.54\textwidth}
\includegraphics[width=\textwidth]{\outputdir/fig/coal_summary_plot.png}
\end{column}
\end{columns}
\end{frame}

\begin{frame}{Where coal production moved, 1985--2005}
\presfig[0.78\textheight]{proportionatecircleprod_diff_allminehuc12_1985_2005.png}
\end{frame}

% ---------------------------------------------------------------- 20
\begin{frame}{IV design: the ARP as a shifter of where coal is mined}
First stage:
\begin{equation*}
\begin{split}
\text{NumMines}_{ct} &= \alpha_0 + \alpha_1\text{NumFacility}_{ct} + \alpha_2\text{UpstreamCoalSulfur\%}_c \cdot \text{Post95}_t \\
&\quad + State_s \cdot Y_t + C_c + \nu_{ct}
\end{split}
\end{equation*}
Second stage:
\begin{equation*}
\text{Violation}_{vct} = \beta_0 + \beta_1\text{NumFacility}_{ct} + \gamma\widehat{\text{NumMines}}_{ct} + State_s \cdot Y_t + C_c + \epsilon_{ct}
\end{equation*}
\begin{itemize}\small
    \item $\text{Post95}_t = \mathbf{1}\{t \geq 1995\}$. Sulfur is time-invariant, so $C_c$ absorbs its level.
    \item Outcome: 1 if utility $c$ had a violation of type $v$ at any point in year $t$. Coefficients are in percentage points.
    \item SEs clustered at the utility level.
\end{itemize}
\end{frame}

% ---------------------------------------------------------------- 23
\begin{frame}{Instrument validity}
\begin{itemize}\small
    \item \textbf{Independence.} The ARP is a national clean-air law set by national energy and emissions politics, not downstream utilities' water quality.
    \item \textbf{Exclusion.} Effect of instrumenting mining downstream of utility on utility behavior is not significant. \backuplink{kgridmrup}{MR violations}, \backuplink{kgridvisitup}{visits}, \backuplink{kgridenfup}{enforcement}.

    Utility characteristics are not predicted by instrument. \backuplink{exclfacilities}{full balance table}
    \item \textbf{Monotonicity.} A defier is a high-sulfur watershed that expanded \emph{because} its coal was high sulfur. Scrubber adoption costs more likely lead to lost market share rather than gain; low-cost high-sulfur regions that expand anyway are never-takers. 
    %Monotonicity is pairwise in sulfur, so low- and high-sulfur areas may respond in opposite directions (\cite{angrist1994identification}).
\end{itemize}
\end{frame}

% ---------------------------------------------------------------- 24
\begin{frame}{Mining makes utilities stop reporting}
\presreg{2sls_dwnstrm_minevio_mr_ivsum_binvio_k2}
\takeaway{Base rates of nitrates, arsenic, and IOC are 5.7\%, 3.7\%, and 4.3\%. OLS estimates are far smaller --- consistent with negative omitted variable bias, as expected if poorly performing utilities (more MR violations) are protected from mine siting (fewer mines). \backuplink{kgridmrdown}{effect across choice of upstream watershed radius}}
\end{frame}

% ---------------------------------------------------------------- 26
\begin{frame}{Reporting cost, not concealment}
\begin{itemize}\small
    \item Nitrate rules require \textbf{more frequent monitoring} once a reading exceeds 50\% of the MCL --- a discrete jump in $c$, with no jump in the penalty for under-reporting and no MCL.
    \item Exceedance risk is held fixed (Channel B = 0), reporting cost increases (Channel A $\geq 0$).
\end{itemize}
\presreg{mr_concentration_lag_ols_k2}
\takeaway{Descriptive, not causal --- read the sign, not the magnitude. Utilities under-report when self-reporting gets expensive, not to conceal exceedances.}
\end{frame}

% ---------------------------------------------------------------- 28
\begin{frame}{Enforcement does not counteract under-reporting}
\presreg{h3_inf_formal_d12_k2}
\takeaway{One more upstream mine reduces formal enforcement by 3.4 pp with utility and year fixed effects alone (significant), but the effect loses significance once state $\times$ year fixed effects are added. \backuplink{kgridenf}{effect across choice of upstream watershed radius}}
\end{frame}

% ---------------------------------------------------------------- 29
\begin{frame}{The regulator increases inspections, not enforcement}
\presreg{h2_snsv_d12_k2}
\takeaway{Inspection visits rise significantly ($+$2.1 pp). Regulators do not systematically step up the visit types tied to sanctions. Visits would reveal concealment and be associated with higher penalties. \backuplink{kgridvisit}{effect across choice of upstream watershed radius}}
\end{frame}

% ---------------------------------------------------------------- 30
\begin{frame}{Conclusion}
\begin{itemize}\small
    \item Exogenous contamination raises the cost of self-reporting, and self-reporting falls --- \textbf{even when utilities remain below MCL}.
    \item Under-reporting due to cost avoidance, rather than concealing a contamination violation.
    \item Regulators' inspection visits rise significantly --- substituting for utility self-reporting.
    \item Enforcement response is mixed and fixed-effects-sensitive: formal action falls with baseline fixed effects but loses significance once state $\times$ year effects are added.
    \item \textbf{Policy.} Tightening MCLs does nothing if nobody tests. Raising self-reporting penalties could reduce under-reporting.
    %\item \textbf{Policy.} Tightening MCLs does nothing if nobody tests. By Proposition 3 the binding lever is the price of \emph{not} reporting where source water is impaired, and it is not currently being pulled (\cite{malik1993self}; \cite{harford1987self}).
    \item Open question: whether regulator monitoring is more efficient than self-reporting.
\end{itemize}
\vfill
\centering \small ek559@cornell.edu
\end{frame}

% ================================================================
% Appendix
% ================================================================
\begin{frame}
\vfill
\centering
\begin{beamercolorbox}[sep=8pt,center,shadow=true,rounded=true]{title}
\usebeamerfont{title}Appendix\par
\end{beamercolorbox}
\vfill
\end{frame}

\begin{frame}{Appendix: combined MR and MCL violations}
\presreg{2sls_dwnstrm_minevio_allcat_ivsum_binvio_k2}
\takeaway{The combined effect mirrors the MR effect: nitrate 3.9--5.1 pp, arsenic 4.0--4.9 pp, IOC 3.1--4.6 pp across the two fixed-effects specifications.}
\end{frame}

\begin{frame}{Appendix: instrument balance on intake facilities}
\hypertarget{exclfacilities}{}
\presreg{exclusion_test_num_facilities_k2}
\end{frame}

\begin{frame}{Appendix: balance test on utility characteristics}
\hypertarget{balancetest}{}
\presreg[max width=\textwidth, max totalheight=0.68\textheight]{pt_balance_6yr_k2}
\takeaway{Observable characteristics jointly predict continuous upstream production only without state fixed effects (p=0.029) and lose significance once state fixed effects are added (p=0.461), as in the main specification. The joint test on the binary presence of upstream mining stays significant with or without state fixed effects.}
\end{frame}

\begin{frame}{Appendix: first stage across choice of upstream watershed radius}
\hypertarget{kgridfs}{}
\presfig[0.62\textheight]{kgrid_first_stage.png}
\takeaway{The first-stage coefficient and F-statistic for utilities downstream of a mine stay well above conventional weak-instrument thresholds as the upstream watershed radius $k$ varies from one to eight flow steps.}
\end{frame}

\begin{frame}{Appendix: exclusion placebo, MR violations across $k$}
\hypertarget{kgridmrup}{}
\presfig[0.66\textheight]{kgrid_mr_upstream_of_mine_coefs.png}
\takeaway{For utilities upstream of a mine --- where mining cannot reach the intake hydrologically --- OLS estimates on MR violations are positive and significant at every $k$. In the reduced form and 2SLS, nitrate estimates are significant for $k \geq 3$, while arsenic and inorganic chemical estimates are significant only at $k=1$.}
\end{frame}

\begin{frame}{Appendix: exclusion placebo, visits across $k$}
\hypertarget{kgridvisitup}{}
\presfig[0.68\textheight]{kgrid_visit_upstream_of_mine_coefs.png}
\takeaway{For utilities upstream of a mine, inspection visits rise significantly at nearly every $k$, sanitary visits at $k=1$ and $k \geq 4$, and sample collection visits at several radii; technical assistance and enforcement visits are mostly insignificant.}
\end{frame}

\begin{frame}{Appendix: exclusion placebo, enforcement across $k$}
\hypertarget{kgridenfup}{}
\presfig[0.76\textheight]{kgrid_enf_upstream_of_mine_coefs.png}
\takeaway{Enforcement actions show no consistent pattern for utilities upstream of a mine at any choice of $k$, supporting the exclusion restriction.}
\end{frame}

\begin{frame}{Appendix: concentration effects across choice of upstream watershed radius}
\hypertarget{kgridsyr2}{}
\presfig[0.62\textheight]{kgrid_syr2_downstream_of_mine_coefs.png}
\takeaway{Arsenic and barium rise significantly only at $k=1$ and are indistinguishable from zero at larger radii. Nitrate is insignificant at $k=1$, significant from $k=2$ to $k=4$ (peaking at $k=3$), and fades beyond; selenium is significant only at $k=1$ and $k=2$.}
\end{frame}

\begin{frame}{Appendix: MR violation effect across choice of upstream watershed radius}
\hypertarget{kgridmrdown}{}
\presfig[0.58\textheight]{kgrid_mr_downstream_of_mine_coefs.png}
\takeaway{The 2SLS effect on MR violations is largest at $k=1$ and declines smoothly as the upstream radius widens. Nitrate stays positive and significant at every $k$; arsenic and inorganic chemical estimates stay positive but are borderline or insignificant at some radii of $k \geq 3$.}
\end{frame}

\begin{frame}{Appendix: enforcement effect across choice of upstream watershed radius}
\hypertarget{kgridenf}{}
\presfig[0.68\textheight]{kgrid_enf_coefs.png}
\takeaway{The formal-enforcement estimate is not stable in sign or significance across choices of $k$, consistent with the fixed-effects sensitivity found in the main specification.}
\end{frame}

\begin{frame}{Appendix: visit effect across choice of upstream watershed radius}
\hypertarget{kgridvisit}{}
\presfig[0.62\textheight]{kgrid_visit_coefs.png}
\takeaway{The rise in inspection visits downstream of mining is significant at every $k \geq 2$ but not at $k=1$. Sanitary visits remain insignificant throughout, while enforcement visits are significant at $k=1$ and at most radii beyond $k=2$.}
\end{frame}

\begin{frame}{Appendix: who appears in the Six-Year Review}
\pressum{syr2_mr_comparison_k2}
\takeaway{Utilities with SYR2 readings are 7--14 pp more likely to ever incur an MR violation. The concentration sample over-represents the more negligent utilities, so it is not selected toward clean systems.}
\end{frame}

\begin{frame}{Appendix: enforcement and visits in the panel}
\pressum{enforcement_visit_type_panels_k2}
\takeaway{Enforcement responds to violations --- any enforcement rises from 21\% overall to 71\% in MR years --- but sanitary visits are no more likely in MR years than in any other year.}
\end{frame}

\begin{frame}{Appendix: residualized first stage}
\presfig[0.72\textheight]{first_stage_scatter_dwnstrm.png}
\end{frame}

\begin{frame}{Appendix: sample watersheds}
\presfig[0.72\textheight]{map_huc12_main_sample.png}
\end{frame}

\begin{frame}{Appendix: Corollary 1, statement and proof}
\scriptsize
Write the share of types under-reporting as $h(\theta,m) = 1 - G(c^*(\theta);m)$, with cutoff $c^*(\theta) = t - (r-ps)\,q(a_{SR}(\theta))$, so that
\[
\frac{\partial MR}{\partial m} = \underbrace{\int_0^{\bar\theta}\!\frac{\partial h(\theta,m)}{\partial m}\,f(\theta;m)\,d\theta}_{\text{A: cost}} \;+\; \underbrace{\int_0^{\bar\theta}\! h(\theta,m)\,\frac{\partial f(\theta;m)}{\partial m}\,d\theta}_{\text{B: composition}}.
\]
\textbf{Corollary 1.} If $q(a_{SR}(\theta)) \equiv \bar q$ on $(0,\bar\theta)$, then $\partial MR/\partial m = -\,\partial G(\bar c;m)/\partial m \ge 0$, where $\bar c = t-(r-ps)\bar q$.
\smallskip

\textit{Proof.} The cutoff is then type-invariant, $c^*(\theta) \equiv \bar c$, so $h(\theta,m) = 1-G(\bar c;m)$ does not vary with $\theta$ and comes outside both integrals. The limits $0$ and $\bar\theta$ do not depend on $m$, so by Leibniz's rule
\[
\text{B} = \big[1-G(\bar c;m)\big]\,\frac{d}{dm}\!\left[\int_0^{\bar\theta}\! f(\theta;m)\,d\theta\right] = \big[1-G(\bar c;m)\big]\,\frac{d}{dm}[1] = 0,
\]
since $f(\cdot\,;m)$ is a density for every $m$. Channel A collapses to $-\,\partial G(\bar c;m)/\partial m \ge 0$, because contamination shifts reporting costs up. $\square$
\smallskip

\textbf{Why it matters.} Channel B is the concealment margin: it operates only if contamination reallocates mass toward types whose exceedance risk, and hence whose cutoff, differs. Shut it down and any remaining rise in under-reporting is avoidance of the reporting technology itself. The 50\%-of-MCL nitrate step is the empirical counterpart --- monitoring frequency jumps, the sanction schedule does not, and readings are nowhere near the limit.
\end{frame}

\begin{frame}[allowframebreaks]{References}
\printbibliography
\end{frame}

\end{document}
